\subsection{Phân tích và so sánh giữa các mô hình}\label{subsec:phan-tich-so-sanh}

\subsubsection{Phân tích kết quả từng mô hình}

\paragraph{1. Logistic Regression (\texttt{src/model\_logistic.py})}\mbox{}
\begin{itemize}
    \item \textbf{Hiệu năng qua 5-fold CV:} Tại ngưỡng $t = 0.5$, mô hình đạt \texttt{AUPRC} ổn định trong khoảng $0.697 - 0.759$ (trung bình $\approx 0.740$) và \texttt{Recall} rất cao ($0.855 - 0.961$). Tuy nhiên, \texttt{Precision} chỉ dao động $0.049 - 0.064$ và \texttt{F1} đạt $0.093 - 0.119$.
    \item \textbf{Hiệu năng trên tập Test:} Mô hình đạt \texttt{AUPRC} $= 0.680$, \texttt{Recall} $= 0.874$ (phát hiện $83/95$ giao dịch gian lận, bỏ sót $\text{FN} = 12$), \texttt{Precision} $= 0.056$ ($\text{FP} = 1{,}403$ trên $56{,}651$ giao dịch bình thường) và \texttt{F1} $= 0.105$.
    \item \textbf{Ưu điểm:} Tốc độ huấn luyện nhanh, kiến trúc đơn giản dễ diễn giải trọng số và nhạy với lớp thiểu số (\texttt{Recall} cao), đóng vai trò mô hình cơ sở (\textit{baseline}) hiệu quả để sàng lọc gian lận.
    \item \textbf{Nhược điểm:} Do là bộ phân lớp tuyến tính nên không biểu diễn được các ranh giới phi tuyến phức tạp; đồng thời việc dùng \texttt{class\_weight="balanced"} khiến mô hình tạo nhiều cảnh báo giả tại ngưỡng mặc định $t = 0.5$.
\end{itemize}

\paragraph{2. Support Vector Machine -- Kernel RBF (\texttt{src/model\_svm.py})}\mbox{}
\begin{itemize}
    \item \textbf{Hiệu năng qua 5-fold CV:} Mô hình đạt \texttt{AUPRC} từ $0.672$ đến $0.807$ (Fold 4 cao nhất đạt $0.807$), \texttt{Recall} trong khoảng $0.737 - 0.868$, \texttt{Precision} đạt $0.492 - 0.735$ và \texttt{F1} đạt $0.628 - 0.772$.
    \item \textbf{Hiệu năng trên tập Test:} Tại $t = 0.5$, mô hình đạt \texttt{AUPRC} $= 0.719$, \texttt{Recall} $= 0.768$ ($73/95$ ca gian lận, $\text{FN} = 22$), \texttt{Precision} $= 0.745$ (chỉ $\text{FP} = 25$ ca báo nhầm) và \texttt{F1} $= 0.756$.
    \item \textbf{Ưu điểm:} Ranh giới phi tuyến của hạt nhân RBF kết hợp chuẩn hóa \texttt{StandardScaler} và giảm mẫu \texttt{RandomUnderSampler(0.05)} giúp cân bằng tốt \texttt{Precision}--\texttt{Recall} ngay tại $t = 0.5$, giảm hơn $56$ lần số cảnh báo giả so với Logistic Regression.
    \item \textbf{Nhược điểm:} Độ nhạy \texttt{Recall} tại $t = 0.5$ thấp hơn Logistic Regression (sót thêm $10$ ca gian lận), thời gian huấn luyện lâu hơn và kết quả giữa các \textit{fold} biến thiên nhẹ do lấy mẫu ngẫu nhiên lớp đa số.
\end{itemize}

\paragraph{3. Artificial Neural Network -- ANN (\texttt{src/model\_ANN.py})}\mbox{}
\begin{itemize}
    \item \textbf{Hiệu năng qua 5-fold CV:} Tại ngưỡng $t = 0.5$, mô hình đạt \texttt{AUPRC} cao trong khoảng $0.723 - 0.870$ (trung bình $\approx 0.818$, cao nhất ở Fold 5 đạt $0.870$), \texttt{Recall} dao động $0.653 - 0.882$, \texttt{Precision} rất cao từ $0.788$ đến $1.000$ (Fold 1 không có báo nhầm nào, $\text{FP} = 0$) và \texttt{F1} đạt $0.741 - 0.855$.
    \item \textbf{Hiệu năng trên tập Test:} Mô hình đạt \texttt{AUPRC} $= 0.825$ (cao nhất trong 4 mô hình), \texttt{Recall} $= 0.747$ (phát hiện $71/95$ ca gian lận, bỏ sót $\text{FN} = 24$), \texttt{Precision} $= 0.887$ (chỉ $\text{FP} = 9$ ca báo nhầm trên $56{,}651$ giao dịch bình thường, $\text{TN} = 56{,}642$) và \texttt{F1} $= 0.811$.
    \item \textbf{Ưu điểm:} Kiến trúc mạng hai tầng ẩn \texttt{(64, 32)} kết hợp hàm kích hoạt ReLU và cơ chế \texttt{early\_stopping=True} giúp học tốt các đặc trưng phi tuyến phức tạp, đạt năng lực xếp hạng xác suất tổng thể (\texttt{AUPRC}) cao nhất đồng thời duy trì tỷ lệ báo nhầm rất thấp ngay tại ngưỡng mặc định $t = 0.5$.
    \item \textbf{Nhược điểm:} Do huấn luyện trên phân phối gốc mất cân bằng mà không áp dụng trọng số lớp nên độ nhạy \texttt{Recall} tại $t = 0.5$ còn bảo thủ (bỏ sót $24$ ca gian lận); hiệu năng giữa các \textit{fold} có độ dao động nhất định (Fold 3 giảm xuống \texttt{AUPRC} $= 0.723$), ít khả năng diễn giải trực tiếp như mô hình tuyến tính và nhạy với siêu tham số khởi tạo.
\end{itemize}

\paragraph{4. Ensemble Learning -- XGBoost (\texttt{src/model\_ensemble.py})}\mbox{}
\begin{itemize}
    \item \textbf{Hiệu năng qua 5-fold CV:} Mô hình duy trì độ ổn định rất cao qua 5 \textit{folds} với \texttt{AUPRC} từ $0.782$ đến $0.859$ (trung bình $\approx 0.822$, cao nhất ở Fold 5 đạt $0.859$), \texttt{Recall} trong khoảng $0.680 - 0.816$, \texttt{Precision} luôn ở mức rất cao ($0.886 - 0.983$) và \texttt{F1} đạt $0.779 - 0.874$.
    \item \textbf{Hiệu năng trên tập Test:} Tại $t = 0.5$, mô hình đạt \texttt{AUPRC} $= 0.764$, \texttt{Recall} $= 0.726$ ($69/95$ ca gian lận, $\text{FN} = 26$), \texttt{Precision} $= 0.958$ (chỉ $\text{FP} = 3$ ca báo nhầm trên $56{,}651$ giao dịch bình thường, $\text{TN} = 56{,}648$) và \texttt{F1} $= 0.826$ (cao nhất trong 4 mô hình tại $t = 0.5$).
    \item \textbf{Ưu điểm:} Tổ hợp 100 cây quyết định nông (\texttt{max\_depth=3}) kết hợp điều chuẩn L2 và thuật toán \texttt{tree\_method="hist"} cho tốc độ huấn luyện rất nhanh, ít biến thiên giữa các \textit{fold} và gần như triệt tiêu cảnh báo giả ($\text{FP} = 3$, \texttt{Precision} $= 95.8\%$), rất phù hợp với dữ liệu bảng.
    \item \textbf{Nhược điểm:} Do chưa thiết lập trọng số lớp (\texttt{scale\_pos\_weight}) hoặc hạ ngưỡng quyết định nên mô hình có xu hướng thận trọng với lớp thiểu số, dẫn đến \texttt{Recall} trên tập Test ($0.726$) thấp nhất trong 4 mô hình tại $t = 0.5$ và \texttt{AUPRC} trên Test ($0.764$) có độ sụt giảm nhẹ so với trung bình CV ($0.822$).
\end{itemize}

\subsubsection{Phân tích so sánh các mô hình tại ngưỡng $t = 0.5$}

Bảng~\ref{tab:compare_4_models_t05} tổng hợp kết quả đánh giá của bốn mô hình trên tập Test tại ngưỡng mặc định $t = 0.5$ sử dụng bộ độ đo thống nhất từ \texttt{src/metrics.py}.

\begin{table}[H]
\centering
\caption{So sánh hiệu năng 4 mô hình trên tập Test tại ngưỡng $t = 0.5$ (\texttt{src/metrics.py})}
\label{tab:compare_4_models_t05}
\renewcommand{\arraystretch}{1.2}
\begin{tabular}{|l|c|c|c|c|c|c|c|c|}
\hline
\textbf{Mô hình} & \textbf{AUPRC} & \textbf{Recall} & \textbf{Precision} & \textbf{F1} & \textbf{TP} & \textbf{FP} & \textbf{FN} & \textbf{TN} \\
\hline
Logistic Regression & $0.680$ & $\mathbf{0.874}$ & $0.056$ & $0.105$ & $83$ & $1403$ & $12$ & $55248$ \\
\hline
SVM (RBF) & $0.719$ & $0.768$ & $0.745$ & $0.756$ & $73$ & $25$ & $22$ & $56626$ \\
\hline
ANN & $\mathbf{0.825}$ & $0.747$ & $0.887$ & $0.811$ & $71$ & $9$ & $24$ & $56642$ \\
\hline
XGBoost & $0.764$ & $0.726$ & $\mathbf{0.958}$ & $\mathbf{0.826}$ & $69$ & $3$ & $26$ & $56648$ \\
\hline
\end{tabular}
\end{table}

\begin{itemize}
    \item \textbf{Chỉ số \texttt{AUPRC}:} \textbf{ANN} dẫn đầu toàn diện với \texttt{AUPRC} $= 0.825$, theo sau là \textbf{XGBoost} ($0.764$), đều vượt trội rõ rệt so với hai mô hình baseline là SVM ($0.719$) và Logistic Regression ($0.680$). Kết quả này cho thấy khả năng học biểu diễn phi tuyến nhiều tầng của ANN giúp tách biệt và xếp hạng xác suất giữa giao dịch gian lận và bình thường hiệu quả nhất trên toàn bộ dải ngưỡng.
    \item \textbf{Chỉ số \texttt{Recall}:} Ở ngưỡng $t = 0.5$, \textbf{Logistic Regression} đạt độ nhạy cao nhất ($0.874$, bắt được $83/95$ ca) nhờ \texttt{class\_weight="balanced"} phạt nặng lỗi bỏ sót ($\sim 600$ lần), tiếp đến là SVM ($0.768$, $73/95$), ANN ($0.747$, $71/95$) và thấp nhất là XGBoost ($0.726$, $69/95$). Nguyên nhân ANN và XGBoost có \texttt{Recall} bảo thủ hơn tại $t = 0.5$ là vì cả hai được huấn luyện trực tiếp trên phân phối dữ liệu gốc mất cân bằng cực mạnh mà chưa điều chỉnh trọng số lớp hay hạ ngưỡng quyết định.
    \item \textbf{Chỉ số \texttt{Precision}:} \textbf{XGBoost} vượt trội tuyệt đối với \texttt{Precision} $= 0.958$ (chỉ $3$ ca báo nhầm trên $56{,}651$ giao dịch bình thường), đứng thứ hai là \textbf{ANN} ($0.887$, $\text{FP} = 9$), bỏ xa SVM ($0.745$, $\text{FP} = 25$) và Logistic Regression ($0.056$, $\text{FP} = 1{,}403$). Nhờ giữ nguyên phân phối gốc cùng cơ chế điều chuẩn cây nông và dừng sớm, cả XGBoost lẫn ANN đều hạn chế tối đa cảnh báo giả.
    \item \textbf{Chỉ số \texttt{F1}:} Vì \texttt{F1} là trung bình điều hòa giữa \texttt{Precision} và \texttt{Recall}, \textbf{XGBoost} đạt mức cân bằng cao nhất tại ngưỡng mặc định $t = 0.5$ với \texttt{F1} $= 0.826$, nhỉnh hơn \textbf{ANN} ($0.811$) và vượt xa SVM ($0.756$) cũng như Logistic Regression ($0.105$). Xét tổng thể tại $t = 0.5$, XGBoost phù hợp nhất cho yêu cầu chính xác cao và ít làm phiền khách hàng, trong khi ANN sở hữu tiềm năng xếp hạng xác suất (\texttt{AUPRC}) tốt nhất nếu tiếp tục tinh chỉnh ngưỡng quyết định để đẩy cao \texttt{Recall}.
\end{itemize}

\subsubsection{Phân tích Precision thấp của Logistic Regression và tác động khi điều chỉnh ngưỡng $t$}

Tại ngưỡng mặc định $t = 0.5$, \texttt{Precision} của Logistic Regression chỉ đạt $0.056$ ($5.6\%$, tức rất thấp dù \texttt{AUPRC} $= 0.680$ vẫn tốt). Nguyên nhân đến từ hai yếu tố cộng hưởng: (1) \textbf{mất cân bằng nền cực đoan} ($56{,}651$ mẫu bình thường so với $95$ mẫu gian lận trên tập Test) khiến chỉ $2.48\%$ tỷ lệ báo nhầm trên lớp \texttt{Normal} cũng tạo ra tới $\text{FP} = 1{,}403$; và (2) tham số \texttt{class\_weight="balanced"} làm dịch xác suất dự đoán $\hat{p}$ lên rất cao, khiến mốc $t = 0.5$ trở thành ngưỡng cắt quá thấp.

Bảng~\ref{tab:lr_threshold_sweep} trình bày kết quả thực nghiệm khi giữ nguyên mô hình Logistic Regression đã huấn luyện và tăng dần ngưỡng quyết định $t$:

\begin{table}[H]
\centering
\caption{Biến thiên các chỉ số của Logistic Regression trên tập Test khi tăng ngưỡng $t$}
\label{tab:lr_threshold_sweep}
\renewcommand{\arraystretch}{1.2}
\begin{tabular}{|c|c|c|c|c|c|c|c|}
\hline
\textbf{Ngưỡng $t$} & \textbf{AUPRC} & \textbf{Precision} & \textbf{Recall} & \textbf{F1} & \textbf{TP} & \textbf{FP} & \textbf{FN} \\
\hline
$0.500000$ (mặc định) & $0.680$ & $0.056$ & $\mathbf{0.874}$ & $0.105$ & $83$ & $1403$ & $12$ \\
\hline
$0.900000$ & $0.680$ & $0.292$ & $0.821$ & $0.431$ & $78$ & $189$ & $17$ \\
\hline
$0.990000$ & $0.680$ & $0.652$ & $0.789$ & $0.714$ & $75$ & $40$ & $20$ \\
\hline
$0.999000$ & $0.680$ & $0.771$ & $0.779$ & $0.775$ & $74$ & $22$ & $21$ \\
\hline
$0.999900$ & $0.680$ & $0.793$ & $0.768$ & $0.781$ & $73$ & $19$ & $22$ \\
\hline
$0.999999$ (tối ưu OOF) & $0.680$ & $\mathbf{0.878}$ & $0.758$ & $\mathbf{0.814}$ & $72$ & $\mathbf{10}$ & $23$ \\
\hline
\end{tabular}
\end{table}

Từ Bảng~\ref{tab:lr_threshold_sweep}, việc điều chỉnh tăng ngưỡng $t$ mang lại ba ý nghĩa quan trọng:
\begin{enumerate}
    \item \textbf{Khẳng định năng lực thực chất của mô hình:} Chỉ số \texttt{AUPRC} $= 0.680$ giữ nguyên không đổi theo $t$. Khi nâng ngưỡng lên $t = 0.999$ hoặc ngưỡng tối ưu $t = 0.999999$ (xác định qua kiểm định chéo \textit{Out-Of-Fold}), \texttt{Precision} tăng mạnh từ $5.6\%$ lên $77.1\% - 87.8\%$, số cảnh báo giả giảm hơn $140$ lần (từ $1{,}403$ xuống còn $10$ ca), đưa \texttt{F1} lên tới $0.814$. Điều này chứng minh nhược điểm tại $t = 0.5$ hoàn toàn nằm ở điểm cắt xác suất chứ không phải do mô hình học kém.
    \item \textbf{Làm rõ sự đánh đổi Precision--Recall:} Khi tăng $t$ từ $0.5$ lên $0.999999$, \texttt{Precision} tăng vọt nhưng \texttt{Recall} giảm tương ứng từ $0.874$ xuống $0.758$ (bỏ sót thêm $11$ ca gian lận). Ngược lại, từ mốc $0.999999$ hạ xuống $0.5$, để bắt thêm $11$ ca gian lận phải đánh đổi bằng $1{,}393$ ca báo nhầm thêm.
    \item \textbf{Ý nghĩa triển khai thực tế:} Khi áp dụng \texttt{class\_weight="balanced"}, cần tinh chỉnh ngưỡng $t$ trên tập kiểm định chéo thay vì dùng cứng $t = 0.5$. Tùy mục tiêu nghiệp vụ ngân hàng, hệ thống có thể chọn $t \approx 0.999999$ nếu muốn tự động khóa thẻ chính xác cao, hoặc chọn ngưỡng trung gian ($t \approx 0.99$) để yêu cầu xác thực bổ sung (OTP) nhằm vừa giữ \texttt{Recall} cao vừa tránh quá tải cảnh báo giả.
\end{enumerate}
